\scp{Aims of this study}{01-01}
The \it{scope} of this study focuses on the prehistory and classification
of the Austronesian languages of modern eastern Indonesia and Timor-Leste.
More specifically, ranging from Sumbawa in the west,
to Aru and the Bird’s Neck of Indonesian Papua in the east,
and to Taliabu and Mangole Islands in the north.
We exclude Sulawesi in the northwest,
and we exclude the South Halmahera-West New Guinea (SHWNG) languages
in the northeast from our primary target region.
However, we do use data from these regions outside our primary
target area for purposes of comparing and contrasting with patterns found within our target region.
And we also look at data from typologically different Papuan languages
still spoken in the region when it appears that there is contact-induced change
in the Austronesian languages within our target region.

The existence of \it{a linguistic subgroup} that encompasses most or
all of the Austronesian languages within our target region has been assumed by some,
and argued for by others.
We are asking the question, “Is there really a basis for such a subgroup?
Are there exclusively shared innovations in the historical phonology,
the morphosyntax, and the lexicon that define and unite such a high-level subgroup,
as well as distinguish it from other subgroups?”

Regardless of the answer to the question of a higher level subgroup,
we are also asking, “What subgroups are found within our target region?”
and “What exclusively shared innovations define those subgroups?”
Subsequent to identifying the various subgroups, we are also asking,
“What sort of links can we find between these subgroups
in our target region that might tie them together at a higher level?”